  \tau(x)(g)=\phi\bigl((x(gu))_{u\in M}\bigr).
\]
The group $G$ is \emph{surjunctive} if every injective cellular automaton on $A^G$ is surjective, for every such $A$. Gottschalk's conjecture asserts that every group is surjunctive~\cite{Gottschalk1973}. The problem asks whether finite local dependence and translation symmetry preserve the finite-set implication from injectivity to surjectivity on the infinite configuration space.

Gromov developed a general symbolic-algebraic surjunctivity theorem~\cite[Section~7.G$'$]{Gromov1999}; Weiss gave the finite-alphabet sofic-group theorem directly~\cite[Theorem~3.2]{Weiss2000}. Elek and Szab\'o explicitly recorded the implication from surjunctivity to direct finiteness over finite fields~\cite[Section~1]{ElekSzabo2004}. Ceccherini-Silberstein and Coornaert developed the linear theory: their equivalence, quantified over all coefficient fields, relates stable finiteness to injective-to-surjective behavior of linear cellular automata with finite-dimensional vector alphabets~\cite[Theorem~4.2]{CSCCoornaert2006}. Phung proved stable finiteness over every field for surjunctive groups, using algebraic cellular automata~\cite[Theorem~B]{Phung2023}. Ceccherini-Silberstein, Coornaert, and Phung later gave a model-theoretic proof and a monoid extension~\cite[Theorem~1.1]{CSCPhungGroup2025,CSCPhungMonoid2025}. These results connect algebraic finiteness to a dynamical property rather than imposing finite dimensionality.

For the finite field in Theorem~\ref{thm:main}, the short algebra-to-automaton passage can be given directly. The coefficients of $b_{\mathrm{out}}$ define a local rule, the coefficients of $a_{\mathrm{out}}$ give a left inverse on all configurations, and a nonzero coefficient of $b_{\mathrm{out}}a_{\mathrm{out}}-1$ identifies a missing point-indicator configuration. Section~\ref{sec:cellular} proves this standard transfer in full, including the multiplication convention. It gives a negative resolution of Gottschalk's conjecture for the group in Theorem~\ref{thm:main}; it requires no further existence theorem or approximation assumption.

